\documentclass[12pt,a4paper]{article}

% ─── Packages ────────────────────────────────────────────────────────────────
\usepackage[a4paper, margin=1in]{geometry}
\usepackage{times}
\usepackage{titlesec}
\usepackage{titling}
\usepackage{fancyhdr}
\usepackage{booktabs}
\usepackage{longtable}
\usepackage{array}
\usepackage{enumitem}
\usepackage{hyperref}
\usepackage{xcolor}
\usepackage{parskip}
\usepackage{setspace}
\usepackage{graphicx}
\usepackage{microtype}
\usepackage{tocloft}
\usepackage{multirow}

% ─── Colours ─────────────────────────────────────────────────────────────────
\definecolor{primary}{RGB}{30, 64, 175}
\definecolor{lightgray}{RGB}{248, 249, 250}
\definecolor{tablegray}{RGB}{230, 235, 245}

% ─── Hyperref Setup ──────────────────────────────────────────────────────────
\hypersetup{
  colorlinks=true,
  linkcolor=primary,
  urlcolor=primary,
  citecolor=primary,
  pdftitle={SRS - Live Interactive Quiz Platform},
  pdfauthor={},
  bookmarks=true,
  bookmarksopen=true
}

% ─── Section Formatting ──────────────────────────────────────────────────────
\titleformat{\section}
  {\large\bfseries\color{primary}}
  {\thesection}{1em}{}[\titlerule]

\titleformat{\subsection}
  {\normalsize\bfseries\color{primary}}
  {\thesubsection}{1em}{}

\titleformat{\subsubsection}
  {\normalsize\bfseries}
  {\thesubsubsection}{1em}{}

\titlespacing*{\section}{0pt}{18pt}{8pt}
\titlespacing*{\subsection}{0pt}{12pt}{5pt}
\titlespacing*{\subsubsection}{0pt}{10pt}{4pt}

% ─── Header / Footer ─────────────────────────────────────────────────────────
\pagestyle{fancy}
\fancyhf{}
\fancyhead[L]{\small\textcolor{primary}{\textbf{PRACTICAL -- 2}}}
\fancyhead[R]{\small\textcolor{primary}{SRS -- Live Interactive Quiz Platform}}
\fancyfoot[C]{\small\thepage}
\renewcommand{\headrulewidth}{0.4pt}
\renewcommand{\headrule}{\hbox to\headwidth{\color{primary}\leaders\hrule height \headrulewidth\hfill}}

% ─── List Formatting ─────────────────────────────────────────────────────────
\setlist[itemize]{noitemsep, topsep=4pt, leftmargin=1.5em}
\setlist[enumerate]{noitemsep, topsep=4pt, leftmargin=1.5em}

% ─── Table of Contents ───────────────────────────────────────────────────────
\renewcommand{\cftsecfont}{\bfseries\color{primary}}
\renewcommand{\cftsubsecfont}{\color{primary}}
\renewcommand{\cftsecpagefont}{\bfseries\color{primary}}

% ─── Custom Commands ─────────────────────────────────────────────────────────
\newcommand{\frbox}[2]{%
  \vspace{2pt}
  \noindent\textbf{#1}\quad #2\par
  \vspace{2pt}
}

% ─────────────────────────────────────────────────────────────────────────────
%  DOCUMENT
% ─────────────────────────────────────────────────────────────────────────────
\begin{document}

% ════════════════════════════════════════════════════════════════════════════
%  TITLE PAGE
% ════════════════════════════════════════════════════════════════════════════
\begin{titlepage}
  \centering
  \vspace*{1.5cm}

  {\LARGE\bfseries\color{primary} PRACTICAL -- 2 \par}
  \vspace{0.4cm}
  \rule{0.8\linewidth}{1.5pt}
  \vspace{0.4cm}

  {\Large\bfseries Software Requirements Specification\\[4pt]
   (SRS) Document Preparation\par}
  \vspace{0.3cm}
  {\large\itshape IEEE Standard\par}

  \vspace{1cm}
  \rule{0.6\linewidth}{0.5pt}
  \vspace{1cm}

  {\large\bfseries Project Title:\par}
  \vspace{0.3cm}
  {\LARGE\bfseries\color{primary} Live Interactive Quiz Platform\par}
  \vspace{0.2cm}
  {\large\itshape (Mentimeter-like Real-Time Quiz System)\par}

  \vspace{1.5cm}
  \rule{0.6\linewidth}{0.5pt}
  \vspace{1cm}

  \begin{tabular}{ll}
    \textbf{Document Type}   & Software Requirements Specification \\[4pt]
    \textbf{Standard}        & IEEE Std 830-1998                   \\[4pt]
    \textbf{Version}         & 1.0                                 \\[4pt]
    \textbf{Date}            & February 2026                       \\[4pt]
    \textbf{Status}          & Final Draft                         \\
  \end{tabular}

  \vfill
  \rule{0.8\linewidth}{1pt}
  {\small Prepared as part of Software Engineering Laboratory Practical}
\end{titlepage}

% ════════════════════════════════════════════════════════════════════════════
%  OBJECTIVE & AIM
% ════════════════════════════════════════════════════════════════════════════
\newpage
\section*{Objective}
To understand the IEEE standard Software Requirements Specification (SRS) format and
prepare a complete SRS document for a software project.

\section*{Aim of the Practical}
The aim of this practical is to train learners to systematically convert a problem
statement into a well-structured SRS document that can serve as a foundation for
software design, development, and testing.

\section*{Introduction}
A Software Requirements Specification (SRS) is a formal document that clearly
describes the functional and non-functional requirements of a software system.
It specifies what the system should do, the constraints under which it must operate,
and the assumptions made during development. The SRS does \emph{not} describe
implementation or design details. A well-written SRS helps avoid ambiguity and
ensures that all stakeholders have a common understanding of the system.

\subsection*{Prerequisite}
Approved Problem Statement (Practical--1) --
\textbf{Live Interactive Quiz Platform (Mentimeter-like)}

% ════════════════════════════════════════════════════════════════════════════
%  TABLE OF CONTENTS
% ════════════════════════════════════════════════════════════════════════════
\newpage
\tableofcontents
\newpage

% ════════════════════════════════════════════════════════════════════════════
%  BODY
% ════════════════════════════════════════════════════════════════════════════
\setcounter{section}{0}

% ─────────────────────────────────────────────────────────────────────────────
\section{Introduction}
% ─────────────────────────────────────────────────────────────────────────────

\subsection{Purpose}
This Software Requirements Specification (SRS) document describes in detail the
functional and non-functional requirements of the \textbf{Live Interactive Quiz
Platform} --- a real-time, Mentimeter-like quiz hosting system built on a
Node.js monorepo architecture with a PostgreSQL database and WebSocket-based
real-time communication.

This document serves as a common reference for quiz creators, developers, testers,
and stakeholders. It clearly defines system behaviour, constraints, and assumptions
so that the system can be developed, tested, and evaluated according to the
specified requirements.

\subsection{Scope}
The \textbf{Live Interactive Quiz Platform} is designed to conduct real-time, live
quiz sessions where registered \emph{creators} host quiz presentations and
registered \emph{participants} respond in real time using a unique 8-digit join code
--- similar to Mentimeter and Kahoot.

\noindent The system will:
\begin{itemize}
  \item Allow registered creators to create and manage quizzes manually, via AI
        prompt generation, or by uploading and scanning a quiz image.
  \item Allow registered participants to join live sessions using a unique 8-digit
        code and answer questions in real time.
  \item Support optional per-question time limits with automatic scoring upon answer
        submission.
  \item Provide a live waiting room, real-time question progression controlled by
        the host, and a live leaderboard updated after each question.
  \item Support adaptive difficulty: automatically adjust the next question's
        difficulty based on the correction rate of the current question
        (optional, creator-enabled).
  \item Provide session result analytics, per-participant response visibility for
        the creator, and export of results to a spreadsheet (CSV).
\end{itemize}

\noindent The scope of the system is limited to live quiz session management and
real-time participation. It does not include payment systems, advanced proctoring,
or integration with external LMS platforms. Session data is automatically retained
for 7--8~days after session expiry and then purged.

\subsection{Definitions, Acronyms, and Abbreviations}

\begin{longtable}{>{\bfseries}p{2.8cm} p{11cm}}
  \toprule
  \textbf{Term} & \textbf{Description} \\
  \midrule
  \endfirsthead
  \toprule
  \textbf{Term} & \textbf{Description} \\
  \midrule
  \endhead
  \bottomrule
  \endfoot
  LIQP  & Live Interactive Quiz Platform \\[3pt]
  Creator & Registered user who creates, manages, and hosts live quiz sessions \\[3pt]
  Participant & Registered user who joins and responds to live quiz sessions \\[3pt]
  AI    & Artificial Intelligence --- used for AI-based quiz generation from text prompts \\[3pt]
  WS / WSS & WebSocket / Secure WebSocket --- real-time bidirectional communication protocol \\[3pt]
  MCQ   & Multiple Choice Question \\[3pt]
  OCR   & Optical Character Recognition --- used to parse quiz content from scanned images \\[3pt]
  LLM   & Large Language Model --- AI model used for generating quiz content from prompts \\[3pt]
  ORM   & Object-Relational Mapping --- Prisma ORM used as the database access layer \\
\end{longtable}

% ─────────────────────────────────────────────────────────────────────────────
\section{Overall Description}
% ─────────────────────────────────────────────────────────────────────────────

\subsection{Product Perspective}
The Live Interactive Quiz Platform is a new standalone web-based application
developed to enable real-time, live quiz hosting and participation, inspired by
platforms like Mentimeter and Kahoot. It is built on a Node.js Turborepo monorepo
with a dedicated HTTP REST backend, a WebSocket real-time backend, and a PostgreSQL
database managed via Prisma ORM.

Currently, most quiz systems are asynchronous and lack live interaction between host
and participants. The proposed system provides a real-time quiz experience where a
host presents questions slide-by-slide to all participants simultaneously, with live
response collection, instant scoring, leaderboard updates after each question, and
detailed per-participant analytics and export for the host.

\subsection{Product Functions}
The Live Interactive Quiz Platform provides the following major functions:

\begin{itemize}
  \item Registered user (creator and participant) authentication via email and password.
  \item Quiz creation in three modes: manual entry, AI prompt-based generation (LLM),
        and image upload with OCR scanning.
  \item Five question types: Single-select MCQ, Multi-select MCQ, True/False, Rating
        Scale, Open-Ended text.
  \item Per-question configuration: time limit (optional), base points, difficulty
        level (1--10).
  \item Live session hosting with a unique 8-digit join code, waiting room lobby,
        and slide-by-slide question control by the host.
  \item Real-time participant response collection and broadcasting via WebSocket.
  \item Automatic scoring: correct answer = full points; incorrect answer = 0 points.
  \item Live leaderboard displayed after each question and at session end.
  \item Adaptive difficulty: optional auto-adjustment of next question difficulty
        based on current question correction rate.
  \item Per-participant response analytics visible to the creator; CSV export of
        session results.
  \item Automatic session data expiry and cleanup after 7--8~days.
\end{itemize}

\subsection{User Classes and Characteristics}

\subsubsection*{Creator}
The creator is a registered user responsible for creating quizzes (manually, via AI
prompt, or by image scan), configuring question types and settings, launching live
sessions, controlling question progression slide-by-slide, pausing/resuming sessions,
and viewing per-participant responses and analytics. The creator can also export
results as a CSV file. The creator is expected to have moderate computer knowledge.

\subsubsection*{Participant}
Participants are registered users who join live quiz sessions using the 8-digit join
code. They answer questions in real time, see their score update after each question,
and view the leaderboard. Late joiners and rejoining after disconnection are both
supported. Participants are assumed to have basic web browsing knowledge and a stable
internet connection.

\subsection{Operating Environment}
The Live Interactive Quiz Platform will operate on:
\begin{itemize}
  \item Desktop computers and smartphones.
  \item Operating systems: Windows, Linux, macOS, and Android.
  \item Modern web browsers supporting WebSocket connections: Chrome, Firefox,
        Edge, Safari.
  \item PostgreSQL relational database for all persistent data storage
        (managed via Prisma ORM).
  \item Node.js runtime environment for HTTP backend and WebSocket backend servers.
  \item A stable internet connection with WebSocket (WS/WSS) protocol support.
  \item External LLM API (e.g., OpenAI GPT) for AI-based quiz generation.
  \item OCR/Vision API for scanning and parsing quiz content from uploaded images.
\end{itemize}

\subsection{Design and Implementation Constraints}
\begin{itemize}
  \item The system shall support the English language interface.
  \item The system shall operate through modern web browsers supporting WebSocket
        connections.
  \item The system shall use WebSocket connections for real-time question
        broadcasting and response collection.
  \item The system shall depend on an external LLM API service (e.g., OpenAI) for
        AI-based quiz generation.
  \item The system shall use Prisma ORM for all database schema management and
        query operations.
  \item Session data shall be retained for a maximum of 7--8~days after session
        expiry.
\end{itemize}

\subsection{Assumptions and Dependencies}
\begin{itemize}
  \item It is assumed that all creators and participants have registered accounts
        with valid email and password credentials.
  \item It is assumed that participants have a stable internet connection during live
        session participation to maintain WebSocket connectivity.
  \item The system depends on the proper functioning of the PostgreSQL database
        server.
  \item The system depends on the availability of the external LLM API (e.g.,
        OpenAI) for AI quiz generation features.
  \item The system depends on WebSocket server availability for real-time session
        management.
  \item The system depends on a file/cloud storage service for quiz image uploads
        used in OCR-based quiz scanning.
\end{itemize}

% ─────────────────────────────────────────────────────────────────────────────
\section{External Interface Requirements}
% ─────────────────────────────────────────────────────────────────────────────

\subsection{User Interfaces}
The system shall provide a simple and user-friendly interface including:
\begin{itemize}
  \item Registration and login page for creators and participants.
  \item Creator dashboard: view all quizzes, create new quiz (manual / AI prompt /
        image scan), manage existing quizzes.
  \item Quiz builder screen: add/edit questions, set question type, time limit,
        points, and difficulty.
  \item Session launch screen: generate 8-digit join code, configure session
        settings (late joiners, rejoin, adaptive difficulty).
  \item Waiting room (lobby) screen: displays join code, shows joined participants,
        creator starts session from here.
  \item Live question screen (participant): displays current question, answer
        options, and countdown timer if timed.
  \item Session control panel (host): advance to next question, pause/resume, end
        session, view live response count.
  \item Leaderboard screen: displayed after each question and at session end.
  \item Analytics screen (creator): view per-question statistics, individual
        participant responses, and export results as CSV.
\end{itemize}

\subsection{Hardware Interfaces}
The system shall operate on standard computers and smartphones with:
\begin{itemize}
  \item Standard keyboard, mouse, or touchscreen for answer input and quiz management.
  \item Display screen for rendering real-time question slides and leaderboard updates.
  \item Network interface card (NIC) supporting stable TCP connections for WebSocket
        communication.
\end{itemize}

\subsection{Software Interfaces}
The system shall use:
\begin{itemize}
  \item \textbf{PostgreSQL} --- relational database for all persistent data (users,
        quizzes, sessions, responses, results).
  \item \textbf{Prisma ORM} --- database access layer with schema migrations and
        type-safe query builder.
  \item \textbf{LLM API} (e.g., OpenAI GPT-4) --- for AI-based quiz generation
        from natural language prompts.
  \item \textbf{OCR / Vision API} --- for scanning and parsing quiz content from
        uploaded images.
  \item \textbf{HTTP Backend} (\texttt{http-backend}) --- Node.js REST API server
        handling user auth, quiz CRUD, session management, and CSV export.
  \item \textbf{WebSocket Backend} (\texttt{ws-backend}) --- Node.js WebSocket
        server for real-time event broadcasting (question transitions, responses,
        leaderboard updates).
  \item \textbf{Frontend} --- Web application for creator and participant interfaces.
\end{itemize}

\subsection{Communication Interfaces}
The system shall use \textbf{HTTP/HTTPS} protocols for all REST API communication
between the client and the HTTP backend server (user authentication, quiz management,
session creation, analytics, CSV export).

The system shall use \textbf{WebSocket (WS/WSS)} protocol for real-time
bidirectional communication between the client and the WebSocket backend server
during live quiz sessions. WebSocket events shall handle: participant joins/leaves,
question broadcasting, response submission acknowledgement, leaderboard updates,
session state changes (start, pause, resume, end), and live participant count updates.

% ─────────────────────────────────────────────────────────────────────────────
\section{System Features (Functional Requirements)}
% ─────────────────────────────────────────────────────────────────────────────

% ── Feature 1 ────────────────────────────────────────────────────────────────
\subsection{Feature 1: User Management}
\textbf{Description:} Manages registration and secure authentication for all users
(creators and participants).

\textbf{Requirements:}
\begin{itemize}
  \item[\textbf{FR1:}] The system shall allow any person to register an account by
        providing a name, email address, and password.
  \item[\textbf{FR2:}] The system shall store passwords as bcrypt-hashed values and
        never in plain text.
  \item[\textbf{FR3:}] The system shall authenticate users via email and password,
        issuing a session token upon successful login.
  \item[\textbf{FR4:}] The system shall maintain a user role field (\texttt{USER} /
        \texttt{ADMIN}) to differentiate access permissions.
  \item[\textbf{FR5:}] The system shall associate each quiz and each hosted session
        with the creator's user account.
  \item[\textbf{FR6:}] The system shall associate each session participation and each
        response with the participant's user account.
\end{itemize}

% ── Feature 2 ────────────────────────────────────────────────────────────────
\subsection{Feature 2: Quiz Creation and Management}
\textbf{Description:} Allows creators to create, configure, edit, and delete quizzes
through three creation modes.

\textbf{Requirements:}
\begin{itemize}
  \item[\textbf{FR7:}] The system shall allow creators to create a quiz manually by
        entering a title, description, cover image (optional), and adding questions
        one by one.
  \item[\textbf{FR8:}] The system shall allow creators to generate a complete quiz by
        entering a natural language AI prompt; the system shall call an external LLM
        API to generate structured quiz content automatically.
  \item[\textbf{FR9:}] The system shall allow creators to upload an image of a quiz;
        the system shall use an OCR/Vision API to parse the image and convert it into
        a structured quiz format.
  \item[\textbf{FR10:}] The system shall allow creators to add questions of the
        following types: Single-select MCQ, Multi-select MCQ, True/False, Rating
        Scale (1--5 or 1--10), and Open-Ended text.
  \item[\textbf{FR11:}] The system shall allow creators to attach an optional image
        to any question.
  \item[\textbf{FR12:}] The system shall allow creators to set a per-question time
        limit (optional checkbox) and specify the time in seconds; if enabled, a
        countdown timer is shown during the live session.
  \item[\textbf{FR13:}] The system shall allow creators to set a base point value for
        each question (default: 100 points).
  \item[\textbf{FR14:}] The system shall allow creators to set a difficulty level
        (1--10) for each question to support adaptive difficulty ordering.
  \item[\textbf{FR15:}] The system shall allow creators to enable or disable the
        Adaptive Difficulty checkbox per quiz.
  \item[\textbf{FR16:}] The system shall allow creators to edit or delete existing
        quizzes and their questions.
  \item[\textbf{FR17:}] The system shall support an optional expiry date on quizzes
        for automatic data cleanup after 7--8~days.
\end{itemize}

% ── Feature 3 ────────────────────────────────────────────────────────────────
\subsection{Feature 3: Live Session Management}
\textbf{Description:} Manages the full lifecycle of a live quiz session from creation
through completion.

\textbf{Requirements:}
\begin{itemize}
  \item[\textbf{FR18:}] The system shall allow a creator to launch a live session
        from any saved quiz, generating a unique 8-digit numeric join code.
  \item[\textbf{FR19:}] The system shall display a waiting room (lobby) after session
        launch where participants can join before the session starts.
  \item[\textbf{FR20:}] The system shall allow registered participants to join a live
        session by entering the 8-digit join code.
  \item[\textbf{FR21:}] The system shall allow late joiners to enter an active session
        if the creator has enabled the \emph{Allow Late Joiners} setting.
  \item[\textbf{FR22:}] The system shall allow a disconnected participant to rejoin an
        active session without losing previously earned points if the creator has
        enabled the \emph{Allow Rejoin} setting.
  \item[\textbf{FR23:}] The system shall allow the host to start the session from the
        waiting room, advancing to the first question and broadcasting it to all
        participants simultaneously via WebSocket.
  \item[\textbf{FR24:}] The system shall allow the host to manually advance to the
        next question (slide-by-slide control); response submission shall be closed
        for the current question before the next one is displayed.
  \item[\textbf{FR25:}] The system shall allow the host to pause and resume the
        session at any time.
  \item[\textbf{FR26:}] The system shall allow the host to end the session, triggering
        display of the final leaderboard.
  \item[\textbf{FR27:}] The system shall automatically expire (invalidate) the
        8-digit join code and schedule session data for cleanup 7--8~days after the
        session ends.
  \item[\textbf{FR28:}] The system shall track the session lifecycle through statuses:
        \texttt{DRAFT} $\rightarrow$ \texttt{WAITING\_ROOM} $\rightarrow$
        \texttt{LIVE} $\rightarrow$ \texttt{PAUSED} $\rightarrow$
        \texttt{COMPLETED}.
\end{itemize}

% ── Feature 4 ────────────────────────────────────────────────────────────────
\subsection{Feature 4: Real-Time Question Interaction}
\textbf{Description:} Handles real-time broadcasting of questions to participants
and collection of their responses during a live session.

\textbf{Requirements:}
\begin{itemize}
  \item[\textbf{FR29:}] The system shall broadcast the current active question to all
        connected participants simultaneously via WebSocket upon host advancement.
  \item[\textbf{FR30:}] The system shall render answer input appropriate to the
        question type: option buttons for MCQ/True-False, a numeric scale for Rating
        Scale, and a text field for Open-Ended.
  \item[\textbf{FR31:}] The system shall display a countdown timer to participants
        for time-limited questions and auto-close response submission when the timer
        reaches zero.
  \item[\textbf{FR32:}] The system shall prevent a participant from modifying or
        resubmitting an answer after it has been submitted for a given question.
  \item[\textbf{FR33:}] The system shall display a live response count to the host
        showing how many participants have answered the current question.
  \item[\textbf{FR34:}] The system shall display the correct answer(s) to all
        participants after response submission is closed for a question.
\end{itemize}

% ── Feature 5 ────────────────────────────────────────────────────────────────
\subsection{Feature 5: Scoring and Leaderboard}
\textbf{Description:} Manages point calculation per response and maintains a
real-time leaderboard across the session.

\textbf{Requirements:}
\begin{itemize}
  \item[\textbf{FR35:}] The system shall award the full question point value to a
        participant who submits a correct answer.
  \item[\textbf{FR36:}] The system shall award zero points to a participant who
        submits an incorrect answer or does not answer within the time limit.
  \item[\textbf{FR37:}] The system shall accumulate total scores for each participant
        across all questions answered in the session, stored in the
        \texttt{SessionParticipant} record.
  \item[\textbf{FR38:}] The system shall calculate and update participant leaderboard
        rankings after each question closes based on total accumulated score.
  \item[\textbf{FR39:}] The system shall broadcast the updated leaderboard to all
        participants via WebSocket after each question, displaying rank, name, avatar,
        and total score.
  \item[\textbf{FR40:}] The system shall display the final leaderboard to all
        participants when the host ends the session.
\end{itemize}

% ── Feature 6 ────────────────────────────────────────────────────────────────
\subsection{Feature 6: Adaptive Difficulty}
\textbf{Description:} Optionally adjusts the difficulty of the next question in a
session based on the live correction rate of the current question.

\textbf{Requirements:}
\begin{itemize}
  \item[\textbf{FR41:}] The system shall allow creators to enable the Adaptive
        Difficulty checkbox per quiz during quiz creation or editing.
  \item[\textbf{FR42:}] After each question closes, the system shall calculate the
        correction rate as:
        \[
          \text{Correction Rate} = \frac{\text{Number of Correct Responses}}
                                        {\text{Total Responses}} \times 100\%
        \]
  \item[\textbf{FR43:}] If adaptive difficulty is enabled and the correction rate is
        high (e.g., $\geq 80\%$), the system shall select the next question with a
        higher difficulty level than the current question.
  \item[\textbf{FR44:}] If adaptive difficulty is enabled and the correction rate is
        low (e.g., $\leq 40\%$), the system shall select the next question with a
        lower difficulty level than the current question.
  \item[\textbf{FR45:}] If adaptive difficulty is enabled and the correction rate is
        between the thresholds (40\%--80\%), the system shall select the next question
        with a similar difficulty level.
  \item[\textbf{FR46:}] The system shall store the correction rate for each question
        in each session in the \texttt{QuestionResult} record for analytics purposes.
\end{itemize}

% ── Feature 7 ────────────────────────────────────────────────────────────────
\subsection{Feature 7: Analytics and Export}
\textbf{Description:} Provides per-session and per-question analytics for the creator
and supports CSV export of session results.

\textbf{Requirements:}
\begin{itemize}
  \item[\textbf{FR47:}] The system shall aggregate and store per-question statistics
        for each session: total responses, correct responses, correction rate (\%),
        and response option distribution (MCQ) or response data (Rating/Open-Ended).
  \item[\textbf{FR48:}] The system shall allow the creator to view a session analytics
        dashboard showing: question-by-question statistics, per-participant responses
        for each question, and overall session summary.
  \item[\textbf{FR49:}] The system shall keep individual participant responses private
        from other participants; only the session creator can view per-participant
        response data.
  \item[\textbf{FR50:}] The system shall allow the creator to export full session
        results as a CSV file with the following columns: Participant Name, Question
        Text, Question Type, Response Given, Is Correct, Points Earned, Response
        Time (ms).
  \item[\textbf{FR51:}] The system shall automatically delete all session data
        (responses, results, participant records) 7--8~days after the session expiry
        date to enforce the data retention policy.
\end{itemize}

% ─────────────────────────────────────────────────────────────────────────────
\section{Non-Functional Requirements}
% ─────────────────────────────────────────────────────────────────────────────
Non-functional requirements define \emph{how well} the system performs rather than
\emph{what} the system does.

\subsection{Performance Requirements}
\begin{itemize}
  \item The system shall respond to HTTP REST API requests within 3~seconds under
        normal operating conditions.
  \item The system shall support at least 2--3 concurrent live sessions at MVP scale
        (scalable independently by the development team later).
  \item The system shall support up to 50 concurrent participants per live session
        at MVP scale.
  \item The system shall broadcast question state updates and leaderboard changes to
        all participants within 500~milliseconds via WebSocket.
  \item The system shall evaluate and record each participant's response within
        1~second of submission.
\end{itemize}

\subsection{Usability Requirements}
\begin{itemize}
  \item The system shall provide a simple and intuitive interface for both creators
        and participants.
  \item The system shall allow creators to launch a live session and start presenting
        questions within 3 steps from the dashboard.
  \item The system shall allow participants to join a live session by entering the
        8-digit code within 2 steps from the home screen.
  \item The system shall clearly display the current question number, total questions,
        timer (if applicable), and submission status to participants during a live
        session.
  \item The system shall minimise the number of steps required for a participant to
        submit an answer.
\end{itemize}

\subsection{Security Requirements}
\begin{itemize}
  \item The system shall require email and password authentication for all users
        (both creators and participants).
  \item The system shall store all passwords as bcrypt-hashed values; plain-text
        passwords shall never be stored or transmitted.
  \item The system shall ensure only the authenticated session host can control
        question progression, pause, resume, or end the session.
  \item The system shall ensure participant responses are private; only the session
        creator can view individual participant responses.
  \item The system shall use HTTPS for all REST API communications and WSS (Secure
        WebSocket) for all real-time communications.
  \item The system shall validate and sanitise all user inputs to prevent injection
        attacks.
\end{itemize}

\subsection{Reliability Requirements}
\begin{itemize}
  \item The system shall maintain session state consistently across the HTTP backend
        and WebSocket backend servers.
  \item The system shall persist each participant response to the PostgreSQL database
        immediately upon submission to prevent data loss in case of server restart or
        network interruption.
  \item The system shall handle participant WebSocket disconnections gracefully,
        preserving all previously earned points and allowing rejoin.
  \item The system shall use Prisma ORM transactions where necessary to ensure
        database consistency during concurrent response submissions.
\end{itemize}

\subsection{Availability Requirements}
\begin{itemize}
  \item The system shall be available 24/7 for quiz creation, management, and session
        hosting.
  \item Scheduled maintenance shall occur during low-usage periods with prior
        notification.
  \item Session join codes and all associated session data (responses, results,
        participant records) shall be automatically purged 7--8~days after session
        expiry via a scheduled cleanup job.
\end{itemize}

\subsection{Maintainability Requirements}
\begin{itemize}
  \item The system shall be modular in design, organised as a Turborepo monorepo
        separating the HTTP REST backend (\texttt{http-backend}), the WebSocket
        real-time backend (\texttt{ws-backend}), the shared database layer
        (\texttt{packages/db} with Prisma ORM), and the frontend into independent
        packages.
  \item The system shall use Prisma schema migrations to manage all database schema
        changes in a version-controlled manner.
  \item The system shall allow easy updates and feature additions to individual
        packages without affecting other parts of the monorepo.
\end{itemize}

\subsection{Portability Requirements}
\begin{itemize}
  \item The system shall be accessible from any device (desktop or mobile) with a
        modern web browser that supports WebSocket connections.
  \item The backend Node.js services shall be deployable on any Node.js-compatible
        server environment (Windows, Linux, macOS).
  \item The system shall use environment variables for all configuration (database
        URL, API keys, ports) to enable easy deployment across different environments.
\end{itemize}

\subsection{Compliance Requirements}
\begin{itemize}
  \item The system shall comply with data privacy requirements by automatically
        deleting all session data (responses, participant records, results) after the
        defined retention period (7--8~days from session expiry).
  \item The system shall store only the minimum required personal data for user
        accounts: name, email address, and bcrypt-hashed password.
  \item The system shall not share individual participant response data with other
        participants; only the session creator has access to per-participant analytics.
\end{itemize}

% ─────────────────────────────────────────────────────────────────────────────
\newpage
\appendix
\section*{Appendix A: Database Schema Summary}
\addcontentsline{toc}{section}{Appendix A: Database Schema Summary}

\begin{longtable}{>{\bfseries}p{3.5cm} p{9cm} p{2.5cm}}
  \toprule
  \textbf{Model} & \textbf{Purpose} & \textbf{Key Fields} \\
  \midrule
  \endfirsthead
  \toprule
  \textbf{Model} & \textbf{Purpose} & \textbf{Key Fields} \\
  \midrule
  \endhead
  \bottomrule
  \endfoot
  User              & Stores all registered users (creators \& participants) & email, passwordHash, role \\[4pt]
  Quiz              & Template/blueprint for a quiz; reusable across sessions & creatorId, isAIGenerated, enableAdaptiveDifficulty, expiresAt \\[4pt]
  Question          & Individual question within a quiz & questionType, difficulty, hasTimeLimit, points \\[4pt]
  Option            & Answer choice for MCQ/True-False questions & isCorrect, order \\[4pt]
  Session           & Live quiz instance with unique join code & joinCode, status, allowLateJoiners, expiresAt \\[4pt]
  SessionState      & Real-time state of an active session & currentQuestionIndex, isAcceptingResponses, participantCount \\[4pt]
  SessionParticipant & Links users to sessions with scoring & totalScore, rank, isActive \\[4pt]
  Response          & Individual participant answer per question & responseData (JSON), isCorrect, pointsEarned \\[4pt]
  ResponseOption    & Tracks selected options (multi-select MCQ) & responseId, optionId \\[4pt]
  QuestionResult    & Aggregated per-question analytics per session & correctionRate, totalResponses, resultData (JSON) \\
\end{longtable}

\end{document}
